\chapter{การทดสอบแปลงทุเรียนภาคสนามและการจัดการปูน}
\label{chap:field_validation}

\section{พื้นที่ศึกษาและการเก็บตัวอย่างดินสวนทุเรียนภาคตะวันออก}
การประเมินสมรรถนะภาคสนามของเครื่องมือวัดและแบบจำลองปัญญาประดิษฐ์ ดำเนินการในแปลงปลูกทุเรียนของเกษตรกรจำนวน 30 แปลง ครอบคลุมพื้นที่ปลูกสำคัญใน 2 จังหวัดของภาคตะวันออก ได้แก่
\begin{enumerate}
    \item \textbf{อำเภอมะขาม จังหวัดจันทบุรี (หมู่บ้านหนองอ้อ):} ดินร่วนปนทราย สภาพดินค่อนข้างเป็นกรด (pH เฉลี่ย 4.8 - 5.2) ปลูกทุเรียนพันธุ์หมอนทองอายุ 8-15 ปี
    \item \textbf{อำเภอท่าใหม่ จังหวัดจันทบุรี (บ้านหนองตาลิ่น ตำบลสองพี่น้อง):} ดินลูกรังชุดดินเขากระทิง สภาพดินกรดปานกลาง (pH เฉลี่ย 5.0 - 5.5) มีการจัดการแปลงแบบเกษตรแปลงใหญ่
    \item \textbf{อำเภอเขาสมิง จังหวัดตราด (หมู่บ้านตาละวาย ตำบลประณีต):} พื้นที่ลุ่มน้ำเขาสมิง ดินเหนียวปนทราย สภาพดินเป็นกรดรุนแรง (pH ต่ำกว่า 4.8) มีประวัติการใช้ปุ๋ยเคมีต่อเนื่อง
\end{enumerate}

\section{เกณฑ์การวินิจฉัยและคำนวณปริมาณปูนโดโลไมต์ปรับปรุงดิน}
ทุเรียนเจริญเติบโตและให้ผลผลิตสูงสุดที่ระดับ pH ดินระหว่าง \textbf{5.5 ถึง 6.5} เมื่อตรวจพบว่าดินมีค่า pH ต่ำกว่า 5.5 ระบบผู้เชี่ยวชาญในแอปพลิเคชันจะคำนวณปริมาณปูนโดโลไมต์ ($\text{CaMg(CO}_3)_2$) หรือปูนมาร์ลที่ต้องใส่เพื่อยกระดับ pH ขึ้นสู่ระดับ 6.0 ตามสูตรคำนวณอิงเนื้อดิน ดังแสดงในตารางที่ \ref{tab:lime_rates}

\begin{table}[H]
\caption{เกณฑ์การใส่ปูนโดโลไมต์เพื่อยกระดับ pH ขึ้น 1.0 หน่วยตามชนิดของเนื้อดิน}
\label{tab:lime_rates}
\centering
\begin{tabularx}{\textwidth}{p{3.5cm}cZZ}
\toprule
\textbf{ชนิดเนื้อดิน} & \textbf{อัตราปูน (กก./ไร่)} & \textbf{ความจุแลกเปลี่ยนแคตไอออน (CEC)} & \textbf{ระยะเวลาพักดินก่อนใส่ปุ๋ย} \\
\midrule
ดินทราย (Sand) & 150 -- 200 & ต่ำ ($< 5\text{ meq/100g}$) & 15 วัน \\
ดินร่วนปนทราย (Sandy Loam) & 250 -- 350 & ปานกลาง ($5 - 15\text{ meq/100g}$) & 20 วัน \\
ดินร่วน (Loam) & 350 -- 450 & ปานกลางค่อนข้างสูง & 25 วัน \\
ดินเหนียว (Clay) & 500 -- 650 & สูง ($> 20\text{ meq/100g}$) & 30 วัน \\
\bottomrule
\end{tabularx}
\end{table}

\section{ผลการทดสอบสมมติฐานทางสถิติ}
การทดสอบความใช้ได้ของวิธีตามขั้นตอนที่ 5.2 ของข้อเสนอโครงการ ได้กำหนดการทดสอบสมมติฐานทางสถิติเพื่อพิสูจน์ความเหนือกว่าของแบบจำลองปัญญาประดิษฐ์ ดังนี้
\begin{itemize}
    \item \textbf{สมมติฐานว่าง ($H_0$):} $\text{RMSE}_{AI} \ge \text{RMSE}_{traditional}$ (แบบจำลอง AI ไม่ได้ช่วยลดความคลาดเคลื่อนในการวัด)
    \item \textbf{สมมติฐานทางเลือก ($H_1$):} $\text{RMSE}_{AI} < \text{RMSE}_{traditional}$ (แบบจำลอง AI ช่วยลดความคลาดเคลื่อนได้อย่างมีนัยสำคัญ)
\end{itemize}

\begin{table}[H]
\caption{ผลการวิเคราะห์สถิติทดสอบสมมติฐานเปรียบเทียบระหว่างวิธีดั้งเดิมกับแบบจำลอง AI}
\label{tab:hypothesis_results}
\centering
\begin{tabularx}{\textwidth}{p{5.5cm}ccZ}
\toprule
\textbf{ดัชนีทางสถิติ} & \textbf{วิธีดั้งเดิม (Traditional)} & \textbf{แบบจำลอง AI (SoilpHTxAI)} & \textbf{การเปลี่ยนแปลง} \\
\midrule
ค่าความคลาดเคลื่อนกำลังสองเฉลี่ย (RMSE) & 0.2690 pH & 0.0382 pH & \textbf{ลดลง 85.80\%} \\
ค่าคลาดเคลื่อนสัมบูรณ์เฉลี่ย (MAE) & 0.2245 pH & 0.0315 pH & \textbf{ลดลง 85.97\%} \\
สัมประสิทธิ์การตัดสินใจ ($R^2$) & 0.8842 & 0.9998 & \textbf{เพิ่มขึ้น 13.07\%} \\
สถิติทดสอบที (Paired $t$-statistic) & --- & 18.425 & $p < 0.001$ \\
\bottomrule
\end{tabularx}
\end{table}

\begin{academicbox}{ข้อสรุปการทดสอบสมมติฐาน (Statistical Conclusion)}
จากผลการทดสอบด้วยสถิติ Paired Samples $t$-test พบว่าค่าสถิติทดสอบ $t = 18.425$ ซึ่งมีค่ามากกว่า $t_{critical} = 2.045$ ที่ระดับนัยสำคัญ $\alpha = 0.05$ อย่างมาก ($p$-value $< 0.0001$) \textbf{จึงปฏิเสธสมมติฐานว่าง ($H_0$) และยอมรับสมมติฐานทางเลือก ($H_1$)} ซึ่งเป็นข้อพิสูจน์เชิงประจักษ์ว่าแบบจำลองปัญญาประดิษฐ์สามารถลดความคลาดเคลื่อนในการวัด pH ของดินได้อย่างมีนัยสำคัญทางสถิติ โดยลดความคลาดเคลื่อนลงได้ถึง 85.80\% เมื่อเทียบกับวิธีเนิร์นสต์ดั้งเดิม
\end{academicbox}

\clearpage
